\section{Compiled Laboratories}
\label{sec:scenes}

Two questions decide whether the compiler is useful. Which laboratories can
the agent produce? And do the features that make an instrument usable by a
robot survive from the real device to the admitted USD?

\begin{figure*}[t]
  \centering
  \includegraphics[width=\textwidth]{fig2-gallery.pdf}
  \caption{Four Lab2Sim laboratories beyond those in Figure~\ref{fig:hero}.
  Each panel shows the delivered Isaac station and, inset, the objects the
  task acts on. Top: two eBench rooms, shown as cutaways with a Lift2 at the
  bench, for removing a ground-glass stopper (teaching room, r11) and racking
  a glass rod (analytical-instrumentation room, r10.1). Bottom: two
  bench-level liquid stations, for dropping a stir bar into a beaker (r5) and
  heating a tube in a water bath (r1).}
  \label{fig:gallery}
\end{figure*}

Figure~\ref{fig:gallery} answers the first question with four stations that
share nothing with Figure~\ref{fig:hero} except the compiler. Two are
room-scale eBench packages in different laboratories, a teaching room and an
analytical-instrumentation room, with a Lift2 already at the bench: a
stoppered flask to open, and a glass rod to seat in a rack. Two are
bench-level liquid stations: a stir bar to drop into a filled
beaker, and a test tube to heat in a beaker water bath on a hotplate. With
the Fehling bath, powder weighing, titration, and drying oven of
Figure~\ref{fig:hero}, the agent's deliveries cover opening, unloading,
weighing, racking, inserting, immersing, and turning a valve.

\begin{figure*}[t]
  \centering
  \includegraphics[width=\textwidth]{fig3-real-spec-usd.pdf}
  \caption{Two end-to-end agent traces read left to right. Manufacturer
  evidence is what the agent searched out; the contract is four fields from
  the specification it wrote; the admitted USD is the scene Isaac renders
  after ConvertAsset promotion (the oven asset under Isaac Sim~4.1, the
  titration station under 4.5). (a) IKA OVEN~125: the door hinge and shelves
  carry from the catalog view to the loaded, open oven. (b) Burette stopcock:
  the manufacturer's PTFE stopcock becomes a revolute joint with a grasp
  handle on the graduated tube, mounted on the stand through a fixed joint.}
  \label{fig:real-spec}
\end{figure*}

Figure~\ref{fig:real-spec} answers the second question on two instruments
the agent compiled end to end. For the OVEN~125, the catalog photograph shows
a right-hinged door swung open onto wire shelves; the specification writes
that down as a revolute \texttt{door\_hinge}, a handle collider, and shelves
declared as placement surfaces; the admitted USD, rendered from the matching
three-quarter view, opens the same way with a beaker on the shelf the
specification named. For the burette, the manufacturer's exploded stopcock
becomes a revolute \texttt{stopcock} joint with a grasp region on its handle,
and in the admitted scene that joint is the part a Lift2 turns. In both rows
the middle column is what made the right column checkable.
